
\documentclass[interno]{gpt}

\titolo{Invio candidatura, repository e pianificazione documenti}
\data{2026-10-29}
\orario{18:00}{19:00}
\luogo{Discord}
\redattore{corro}
\verificatore{mazzurana}

\modifica{1.0.0}{2026-10-29}{corro}{mazzurana}{Verifica e approvazione}
\modifica{0.1.0}{2026-10-29}{corro}{mazzurana}{Prima stesura}

\componenti{biasini,cescobolla,corro,guglielmo,mazzurana,spagnoli}
\begin{document}

\section{Ordine del giorno}
\begin{enumerate}
  \item Revisione finale dei documenti relativi alla candidatura.
  \item Invio della lettera di presentazione.
  \item Rifinimento del repository: conformità della nomenclatura e gestione delle versioni.
  \item Pianificazione della stesura del Glossario e del Diario di Bordo.
\end{enumerate}

\clearpage

\section{Svolgimento}
\subsection{Candidatura e Lettera di Presentazione}
Il gruppo si è riunito per effettuare un'ultima revisione congiunta di tutta la documentazione prodotta per la candidatura. Constatata la correttezza del materiale e l'allineamento con i requisiti richiesti, il team ha proceduto ufficialmente all'invio della lettera di presentazione.

\subsection{Rifinimento Repository, Nomenclatura e Versioni}
Si è discusso su come strutturare e rifinire il repository per i prossimi sviluppi. Il gruppo ha concordato la necessità di applicare rigorosamente le regole di conformità per la nomenclatura dei file e di standardizzare il sistema di versionamento. Questo passaggio è ritenuto fondamentale per mantenere il lavoro ordinato, tracciabile e coerente con le Norme di Progetto in via di definizione.

\subsection{Organizzazione e ripartizione dei compiti}
Per garantire il rispetto delle prossime scadenze, sono state assegnate le attività relative alla prima stesura del Glossario (versione 1.0.0) e alla preparazione del Diario di Bordo in vista della scadenza del 2 novembre. I carichi di lavoro sono stati distribuiti in coppie per facilitare la revisione incrociata.

\section{Decisioni}
\begin{decisioni}
  \decisione{Approvazione definitiva dei documenti di candidatura e invio della lettera di presentazione.}
  \decisione{Adozione immediata delle linee guida per la conformità della nomenclatura e il versionamento all'interno del repository.}
\end{decisioni}

\section{Azioni da intraprendere}
\begin{azioni}
  \azione[\#57]{Rifinimento repo, conformità nomenclatura e versioni.}{guglielmo, cescobolla}{31/10/2026}
  \azione[\#58]{Stesura Glossario 1.0.0.}{corro, biasini}{01/11/2026}
  \azione[\#51]{Preparazione diario di bordo per il 2 novembre.}{spagnoli, mazzurana}{01/11/2026}
\end{azioni}

\section{Prossima riunione}
La prossima riunione interna è fissata per \textbf{martedì 3 novembre 2026 alle ore 21:00} sulla piattaforma Discord. L'obiettivo principale dell'incontro sarà la presa visione dell'assegnazione ufficiale degli appalti e la conseguente ripianificazione delle attività.

\end{document}
